\documentclass[12pt,a4paper]{article}

% ============================================================================
% KLAUSUR-TEMPLATE: IT-SICHERHEIT (90 Minuten)
% ============================================================================
% Automatisch generiert: 27.05.2026
% Basierend auf Modulhandbuch, Skripten, Notizen und Übungsaufgaben
% ============================================================================

\usepackage[utf8]{inputenc}
\usepackage[german]{babel}
\usepackage[a4paper, margin=2cm]{geometry}
\usepackage{fancyhdr}
\usepackage{amsmath}
\usepackage{amssymb}
\usepackage{listings}
\usepackage{xcolor}
\usepackage{graphicx}
\usepackage{tikz}
\usepackage{tabularx}
\usepackage{enumerate}

% Code-Styling
\lstset{
    language=Python,
    basicstyle=\ttfamily\small,
    keywordstyle=\color{blue},
    commentstyle=\color{gray},
    stringstyle=\color{red},
    breaklines=true,
    showstringspaces=false
}

% Header/Footer
\pagestyle{fancy}
\fancyhf{}
\rhead{IT-Sicherheit Klausur}
\lhead{Bearbeitungszeit: 90 Minuten}
\cfoot{\thepage}

% ============================================================================
\title{\textbf{Klausur: IT-Sicherheit}}
\author{DHBW Stuttgart -- Informatik}
\date{27.05.2026}

\begin{document}

\maketitle

% ============================================================================
% ANWEISUNGEN
% ============================================================================

\section*{KLAUSURANWEISUNGEN}

\begin{itemize}
    \item \textbf{Bearbeitungszeit:} 90 Minuten
    \item \textbf{Hilfsmittel:} Keine
    \item \textbf{Aufgaben:} Alle Aufgaben sind Pflichtaufgaben
    \item \textbf{Punkteverteilung:} Siehe Aufgabenblatt
    \item \textbf{Bestanden:} 50\% der Gesamtpunkte
\end{itemize}

\textbf{Hinweise:}
\begin{itemize}
    \item Schreiben Sie leserlich
    \item Beantworten Sie die Fragen präzise und begründen Sie Ihre Antworten
    \item Nutzen Sie Diagramme, um komplexe Konzepte zu erklären
\end{itemize}

\newpage

% ============================================================================
% AUFGABEN
% ============================================================================

\section{Grundlagen der IT-Sicherheit [15 Punkte]}

\textbf{Aufgabe 1.1 - Schutzziele (6 Punkte)}

Nennen Sie die sechs Schutzziele der IT-Sicherheit und erläutern Sie kurz, was unter dem Schutzziel \textit{Vertraulichkeit} zu verstehen ist. Geben Sie ein Beispiel, wie dieses Schutzziel in der Praxis verletzt werden kann.

\vspace{4cm}

\textbf{Aufgabe 1.2 - Bedrohungsklassifikation (9 Punkte)}

\begin{enumerate}[(a)]
    \item Klassifizieren Sie folgende Bedrohungen und geben Sie jeweils ein Beispiel an:
    \begin{itemize}
        \item Höhere Gewalt
        \item Fahrlässigkeit
        \item Vorsatz
    \end{itemize}
    
    \item Unterscheiden Sie zwischen \textit{passiven} und \textit{aktiven} Angriffen und nennen Sie je ein Beispiel.
\end{enumerate}

\vspace{5cm}

% ============================================================================

\newpage

\section{Symmetrische Verschlüsselung [20 Punkte]}

\textbf{Aufgabe 2.1 - Caesar-Chiffre (8 Punkte)}

Gegeben ist eine Caesar-Chiffre mit Schlüssel $k = 3$ und Blockgröße 1 (ECB-Modus).
Verwenden Sie die Zuordnung: A $\to$ 0, B $\to$ 1, ..., Z $\to$ 25

\begin{enumerate}[(a)]
    \item Verschlüsseln Sie die Nachricht ``SICHERHEIT''
    \item Welche Sicherheitsprobleme hat die ECB-Modus Verschlüsselung?
    \item Nennen Sie einen sichereren Betriebsmodus und begründen Sie die Verbesserung
\end{enumerate}

\vspace{4cm}

\textbf{Aufgabe 2.2 - AES und DES (12 Punkte)}

\begin{enumerate}[(a)]
    \item Vergleichen Sie DES und AES hinsichtlich:
    \begin{itemize}
        \item Schlüssellänge
        \item Blockgröße
        \item Anzahl der Runden
        \item Sicherheitsniveau (aktuell)
    \end{itemize}
    
    \item Warum ist DES heute nicht mehr sicher?
    
    \item Welche Verschlüsselungsstandards sind aktuell für Industrie 4.0 und IoT empfohlen?
\end{enumerate}

\vspace{5cm}

% ============================================================================

\newpage

\section{Kryptologie und Hashfunktionen [18 Punkte]}

\textbf{Aufgabe 3.1 - Hashfunktionen (9 Punkte)}

\begin{enumerate}[(a)]
    \item Nennen Sie mindestens 3 Anforderungen an sichere Hashfunktionen
    
    \item Eine Hashfunktion $H$ erfüllt folgende Eigenschaft nicht:
    ``Preimage-Resistance (Erste Urbild-Resistenz)''. Welche Angriffe werden dadurch ermöglicht?
    
    \item Erklären Sie, warum SHA-1 (160 Bit) kritisch ist, während SHA-256 noch sicher ist.
\end{enumerate}

\vspace{3cm}

\textbf{Aufgabe 3.2 - Message Authentication Code (9 Punkte)}

\begin{enumerate}[(a)]
    \item Unterscheiden Sie zwischen HMAC und einfachen Keyed-Hash-Verfahren
    
    \item Skizzieren Sie einen Angriff auf ``Secret-Prefix'' und ``Secret-Suffix'' Verfahren
    
    \item Warum ist HMAC-SHA-256 in modernen Protokollen dem einfachen Keyed-Hash vorzuziehen?
\end{enumerate}

\vspace{3cm}

% ============================================================================

\newpage

\section{Asymmetrische Verfahren und RSA [17 Punkte]}

\textbf{Aufgabe 4.1 - RSA Grundlagen (10 Punkte)}

Gegeben: Öffentlicher Schlüssel $(e=5, n=21)$, Privater Schlüssel $(d=17)$

\begin{enumerate}[(a)]
    \item Verschlüsseln Sie die Nachricht ``5'' mit dem öffentlichen Schlüssel
    
    \item Entschlüsseln Sie die Nachricht ``11'', die mit dem öffentlichen Schlüssel verschlüsselt wurde
    
    \item Erklären Sie, warum dieses Beispiel in der Praxis nicht sicher ist
    
    \item Wie groß sollte $n$ mindestens sein, um RSA praktisch sicher zu machen?
\end{enumerate}

\vspace{4cm}

\textbf{Aufgabe 4.2 - Digitale Signatur (7 Punkte)}

\begin{enumerate}[(a)]
    \item Erklären Sie die Funktionsweise einer digitalen Signatur
    
    \item Welche Schutzziele werden durch digitale Signaturen erreicht?
    
    \item Ist es möglich, dass ein Angreifer eine gültige Signatur erzeugt, wenn der private Schlüssel des Signierers geschützt ist?
\end{enumerate}

\vspace{3cm}

% ============================================================================

\newpage

\section{Authentifikation und Identifikation [15 Punkte]}

\textbf{Aufgabe 5.1 - Passwort-Authentifikation (8 Punkte)}

\begin{enumerate}[(a)]
    \item Unterscheiden Sie zwischen Identifikation und Authentifikation
    
    \item Erklären Sie den Unterschied zwischen PAP und CHAP und deren Sicherheitsimplikationen
    
    \item Warum ist das Speichern von Passwort-Hashes mit Salt sicherer als das Speichern von Passwörtern im Klartext?
\end{enumerate}

\vspace{3cm}

\textbf{Aufgabe 5.2 - Mehfaktor-Authentifikation (7 Punkte)}

\begin{enumerate}[(a)]
    \item Nennen Sie die drei Faktoren der Authentifikation
    
    \item Geben Sie zwei Beispiele für Zweifaktor-Authentifikation und erklären Sie deren Vorteile
    
    \item Welche Angriffe können durch MFA nicht verhindert werden?
\end{enumerate}

\vspace{3cm}

% ============================================================================

\newpage

\section{DSGVO und Sicherheits-Engineering [15 Punkte]}

\textbf{Aufgabe 6.1 - Datenschutz-Grundverordnung (10 Punkte)}

\begin{enumerate}[(a)]
    \item Definieren Sie ``personenbezogene Daten'' gemäß DSGVO
    
    \item Nennen Sie mindestens 4 Kategorien von ``Daten besonderer Kategorie''
    
    \item Ein Unternehmen speichert Benutzer-Gesichtsererkennungsdaten. Welche DSGVO-Artikel sind relevant?
    
    \item Welche Strafen kann ein Unternehmen bei DSGVO-Verstößen erhalten?
\end{enumerate}

\vspace{3cm}

\textbf{Aufgabe 6.2 - Security Engineering (5 Punkte)}

\begin{enumerate}[(a)]
    \item Erklären Sie das Konzept ``Secure by Design''
    
    \item Nennen Sie zwei Maßnahmen zur Umsetzung von ``Secure by Default''
\end{enumerate}

\vspace{2cm}

% ============================================================================

\newpage

\section{Praxisaufgabe: Sicherheitsanalyse [10 Punkte]}

\textbf{Aufgabe 7 - Szenarien-Bewertung}

Beurteilen Sie folgende Szenarien unter Sicherheitsgesichtspunkten:

\textbf{Szenario A: Cloud-Authentifikation}

Ein SaaS-Anbieter speichert Benutzerdaten in der Cloud und verschlüsselt diese mit AES-256-CBC. Die Schlüsselverwaltung erfolgt durch den Cloud-Provider. Der Zugriff ist durch Passwörter mit Zwei-Faktor-Authentifikation (SMS-Codes) geschützt.

\begin{enumerate}[(i)]
    \item Welche Schutzziele sind abgedeckt?
    \item Welche Risiken bestehen noch?
    \item Wie würden Sie die Sicherheit verbessern?
\end{enumerate}

\vspace{4cm}

\textbf{Szenario B: IoT-Geräte}

Ein Hersteller von smarten Türschlössern nutzt pre-shared keys für die Authentifikation zwischen Schloss und App. Die Schlüssel sind in die Firmware des Schlosses eingebettet.

\begin{enumerate}[(i)]
    \item Identifizieren Sie die Sicherheitsprobleme
    \item Wie würde ein Angreifer vorgehen?
    \item Welche Lösung würden Sie empfehlen?
\end{enumerate}

\vspace{4cm}

% ============================================================================

\newpage

\section*{PUNKTEVERTEILUNG}

\begin{tabularx}{\textwidth}{|l|X|r|}
    \hline
    \textbf{Aufgabe} & \textbf{Thema} & \textbf{Punkte} \\
    \hline
    1 & Grundlagen IT-Sicherheit & 15 \\
    2 & Symmetrische Verschlüsselung & 20 \\
    3 & Kryptologie und Hashfunktionen & 18 \\
    4 & Asymmetrische Verfahren und RSA & 17 \\
    5 & Authentifikation und Identifikation & 15 \\
    6 & DSGVO und Security Engineering & 15 \\
    7 & Praxisaufgabe: Sicherheitsanalyse & 10 \\
    \hline
    & \textbf{SUMME} & \textbf{110 Punkte} \\
    \hline
\end{tabularx}

\vspace{1cm}

\textbf{Bewertungsskala:}
\begin{itemize}
    \item 1,0: $\geq 99$ Punkte (90\%)
    \item 1,3: $\geq 90$ Punkte (82\%)
    \item 1,7: $\geq 82$ Punkte (75\%)
    \item 2,0: $\geq 77$ Punkte (70\%)
    \item 2,3: $\geq 69$ Punkte (63\%)
    \item 2,7: $\geq 60$ Punkte (55\%)
    \item 3,0: $\geq 55$ Punkte (50\%)
    \item 4,0: $< 55$ Punkte
\end{itemize}

\end{document}
